\documentclass[10pt,letterpaper]{article}
\usepackage[ficha]{../latex_common/estilo_monitoria}
\usepackage{multicol}

\fancyhead[L]{\footnotesize\sffamily\color{GrisTexto}\textbf{Ficha Resumen 01 --- Símbolos y Reglas de DFD}}
\fancyhead[R]{\footnotesize\sffamily\color{Turquesa}\textbf{Monitoría Algoritmos 30h}}

\begin{document}
\thispagestyle{empty}

\begin{tcolorbox}[enhanced,colback=AzulMarino,colframe=AzulMarino,arc=4pt,top=3.5mm,bottom=3.5mm,left=4mm,right=4mm,boxrule=0pt]
  \tcbox[on line,colback=Turquesa,colframe=Turquesa,arc=2pt,boxsep=1pt,left=4pt,right=4pt]{\footnotesize\sffamily\bfseries\color{white} FICHA IMPRIMIBLE 01}\hspace{6pt}%
  {\large\bfseries\sffamily\color{white} Símbolos Oficiales y Reglas de Oro de DFD}\\[2pt]
  {\footnotesize\sffamily\color{AzulSuave} Guía visual rápida para dibujar Diagramas de Flujo sin errores en clase y en parciales}
\end{tcolorbox}

\vspace{2pt}
\begin{center}
  \includegraphics[width=0.98\textwidth,height=9.8cm,keepaspectratio]{dfd_00_simbolos_oficiales.png}
\end{center}

\vspace{-2pt}
\begin{multicols}{2}
\begin{consejo}{Las 6 Reglas de Oro de todo DFD}
\footnotesize
\begin{enumerate}[leftmargin=12pt,itemsep=2pt]
  \item \textbf{Un solo Inicio y un solo Fin:} Todo diagrama comienza en un óvalo \emph{Inicio} y cierra en un óvalo \emph{Fin}.
  \item \textbf{Flechas con sentido único:} Siempre muestran hacia dónde fluye el algoritmo (de arriba hacia abajo). Nunca dejes flechas sueltas.
  \item \textbf{Rombo etiquetado:} Cada decisión tiene exactamente dos salidas claras: \textbf{Sí / V} y \textbf{No / F}.
  \item \textbf{Convergencia limpia:} Tras un \emph{Sí/No}, las ramas vuelven a unirse antes de seguir (o terminan en ciclo).
  \item \textbf{Inicializa antes del ciclo:} Contadores (\codigo{c <- 0}) y acumuladores (\codigo{s <- 0}) van \textbf{fuera y antes} del ciclo.
  \item \textbf{Prueba en papel:} Simula el diagrama con 3 números reales antes de pasarlo a PSeInt o C++.
\end{enumerate}
\end{consejo}

\columnbreak

\begin{pruebaescritorio}{Ejemplo clásico: Suma y Promedio de 3 Notas}
\centering
\includegraphics[width=0.92\linewidth,height=5.2cm,keepaspectratio]{dfd_01_suma_promedio.png}
\vspace{2pt}
\scriptsize
\begin{tabular}{ccccc}
\toprule
\textbf{n1} & \textbf{n2} & \textbf{n3} & \textbf{prom = (n1+n2+n3)/3} & \textbf{Salida} \\
\midrule
4.0 & 3.5 & 4.5 & 12.0 / 3 = 4.0 & \textbf{4.0} \\
\bottomrule
\end{tabular}
\end{pruebaescritorio}
\end{multicols}

\vspace{-4pt}
\begin{peligro}{Errores que anulan puntos en DFD}
\footnotesize
\textbf{1)} Usar un rectángulo para leer del teclado (leer usa paralelogramo/trapecio). \textbf{2)} Olvidar actualizar la variable de control dentro de un ciclo \emph{Mientras} (bucle infinito). \textbf{3)} Escribir código C++ dentro de las cajas en vez de expresiones claras como \codigo{prom <- (n1+n2+n3)/3}.
\end{peligro}

\end{document}
